% GENERATED by gen_docs.py from terms.yaml — do not edit
\begin{description}
\item[{\hypertarget{term:alias-image}{alias image}}] A recording cannot hold tones above half its sample rate; higher ones fold back down and can land on a harmonic being measured.
\item[{\hypertarget{term:analysis-lattice}{analysis lattice}}] The grid of frequency slots the two notes and all their invented tones land on, arranged so that no two invented tones share a slot or crowd each other.
\item[{\hypertarget{term:analysis-window}{analysis window} ($N$)}] The long middle stretch of the recording that the analysis actually reads, cut so that neither the start nor the end of the notes is in it.
\item[{\hypertarget{term:averaged-cycle}{averaged cycle}}] Every cycle of the note laid on top of the previous ones and averaged, so that noise cancels and the shape stays --- the figure is drawn from that one averaged cycle.
\item[{\hypertarget{term:box-normalization}{box normalization}}] Loop areas are stated as a fraction of the box the whole figure fits in, so a tiny signal and a loud one are on the same scale.
\item[{\hypertarget{term:capture-floor}{capture floor}}] How deep this particular recording could see, stated under the table; it describes the recording and never changes a verdict.
\item[{\hypertarget{term:clipping}{clipping}}] What a distortion circuit does to a note that is too big for it: the peaks of the wave are flattened or rounded off. A wave whose shape has changed is, by the same token, a wave with new harmonics in it.
\item[{\hypertarget{term:coherence}{coherence} ($\Gamma$)}] A score for how orderly a harmonic's measurement is across neighbouring frequencies; real circuit behaviour is orderly, noise is not.
\item[{\hypertarget{term:coherence-check}{coherence check}}] A check that neither note drifted while the recording ran; if one did, the whole record is smeared and the app says so rather than quoting numbers.
\item[{\hypertarget{term:coherence-detector}{coherence detector}}] The nine slots either side of each note, kept with the record so that a note which drifted during the capture can be caught where its spill lands first.
\item[{\hypertarget{term:coherent-sampling}{coherent sampling}}] Both notes are tuned so they fit the analysis window exactly, which is what lets each invented tone be read cleanly in its own slot.
\item[{\hypertarget{term:coupling-capacitor}{coupling capacitor}}] The capacitor every pedal has on its output. It removes the wave's average level, so a clipper that clips the top and bottom at different heights shows equal heights after it --- but the even harmonics that asymmetry makes are still there.
\item[{\hypertarget{term:dbc}{dBc}}] A level relative to the note (or tone) that caused it, so that ``how big is the distortion'' does not depend on how loud the test was.
\item[{\hypertarget{term:dbfs}{dBFS}}] A level relative to the loudest a digital recording can be; 0 is the maximum and negative numbers are quieter.
\item[{\hypertarget{term:decibel}{decibel} (dB)}] A way of writing how much bigger or smaller one level is than another, on a scale where each step of 20 is a factor of ten in size: 0 means the same, $-20$ means a tenth as big, $-40$ a hundredth, $-60$ a thousandth.
\item[{\hypertarget{term:deconvolution}{deconvolution}}] The step that undoes the sweep, turning the recording into a set of separate blips, one per harmonic.
\item[{\hypertarget{term:derotation}{derotation}}] Lining up repeated sweeps to the fraction of a sample before combining them, so that a tiny timing drift cannot smear the result.
\item[{\hypertarget{term:excited-band}{excited band}}] The range of notes the test tone actually played at full strength; the very first and very last moments are ramped and do not count.
\item[{\hypertarget{term:extraction-window}{extraction window} ($T_k$)}] The slice of time around one harmonic's blip that is kept for measuring it; wide enough to hold the blip, narrow enough to exclude its neighbours.
\item[{\hypertarget{term:frame-fit}{frame fit}}] Lining the two measurements up in time by one shift that makes every harmonic agree at once; how well they then agree says whether the two records really describe the same device at the same setting.
\item[{\hypertarget{term:fundamental}{fundamental} ($f_0$)}] The note being played. Distortion shows up as extra tones at whole-number multiples of it.
\item[{\hypertarget{term:fundamental-ellipse}{fundamental ellipse}}] The smooth oval a plain tone filter draws by shifting the note in time: real, and not memory, so it is measured separately and subtracted.
\item[{\hypertarget{term:fundamental-presence}{fundamental presence}}] A check that the note itself is still there in the output; if a device removes the note, ratios to it mean nothing and the app says so.
\item[{\hypertarget{term:harmonic-order}{harmonic order} ($k$)}] Which multiple of the note an added tone is: twice the note is the second harmonic, three times is the third, and so on.
\item[{\hypertarget{term:harmonic-pulse}{harmonic pulse}}] After the analysis untangles the recording, each harmonic shows up as its own separate blip, spaced apart in time so it can be measured alone.
\item[{\hypertarget{term:harmonic-response}{harmonic response} ($H_k(f)$)}] For each multiple of the note, how much of it the device adds, across the whole range of notes.
\item[{\hypertarget{term:harmonic-role}{harmonic role}}] Which of five jobs an invented tone does to the chord you played: thickens it, adds weight below it, adds a colour that is in tune with the chord, clashes, or --- all together --- how much there is.
\item[{\hypertarget{term:shaping}{harmonic shaping}}] A check that a measured harmonic is really the circuit's doing and not just the device's hiss showing up where a harmonic would be.
\item[{\hypertarget{term:impulse-response}{impulse response} ($h[n]$)}] The device's fingerprint in time: a sharp blip for the note itself and earlier, smaller blips for each harmonic.
\item[{\hypertarget{term:input-normalized}{input-normalized}}] Measured relative to the strength of the test tone, so that the numbers compare across drive levels and rigs.
\item[{\hypertarget{term:instantaneous-frequency}{instantaneous frequency} ($f(t)$)}] The pitch the test tone is at, at a given moment.
\item[{\hypertarget{term:intermodulation-product}{intermodulation product} ($m f_1 + n f_2$)}] An extra tone the pedal makes from the two notes together --- a sum, a difference, or a combination of multiples of the two --- and the recipe that names which combination it is.
\item[{\hypertarget{term:inverse-filter}{inverse filter} ($\tilde X^{-1}(f)$)}] The exact mathematical undo of the test tone, written down from its formula rather than measured, so it keeps working above the highest note the tone reached.
\item[{\hypertarget{term:just-ratio}{just ratio}}] The whole-number version of the interval you asked for --- three to two for a fifth --- which decides where each invented tone falls musically.
\item[{\hypertarget{term:lattice-snap}{lattice snap}}] The small, disclosed adjustment of the two notes --- a few cents --- that lets the analysis read every invented tone cleanly.
\item[{\hypertarget{term:lattice-spacing}{lattice spacing} ($g$)}] How many slots apart any two invented tones must be; the probe is chosen so this is more than a dozen.
\item[{\hypertarget{term:local-floor}{local floor}}] The noise measured right next to each invented tone, from slots that hold nothing else; a tone counts as present when it stands well above that.
\item[{\hypertarget{term:floor}{measurement floor}}] How quiet a measurement can go before it is measuring the recording equipment instead of the pedal; numbers at that level are stated as limits, not values.
\item[{\hypertarget{term:memory-verdict}{memory verdict}}] The app's three-way call on the loop that is left once the tone filter's oval is taken away: small enough to be treated as no memory, large enough to be called memory, or in between --- where the app hedges, because a tone filter ahead of the clipper draws the same kind of loop as mild memory and one note cannot tell them apart.
\item[{\hypertarget{term:lattice}{memoryless lattice}}] The two positions a harmonic can sit in, relative to the note, in a pedal with no memory and no filtering: exactly in step with it, or exactly reversed. How far each measured harmonic sits from the nearer of the two is what the filters added, and the removal takes that out.
\item[{\hypertarget{term:manifest}{methods manifest}}] The list of every setting the app uses to make a measurement, written out by the app's own code so that it cannot be typed wrongly. The documents read their numbers from it.
\item[{\hypertarget{term:nonlinear-loop-area}{nonlinear loop area}}] The loop that is left once the tone filter's oval is taken away: the part only a real memory effect (or a filter ahead of the clipper, see the phase removal) can draw.
\item[{\hypertarget{term:orientation}{orientation resolution}}] A check that the picture is the right way up: a plugin measurement can come out upside down, and the pedal's own measured polarity is used to turn it back.
\item[{\hypertarget{term:pairing-residual}{pairing residual}}] A single number for how well the two records agree once lined up; too large, and the app declines to apply the correction and says so.
\item[{\hypertarget{term:parity}{parity test}}] A check, run automatically on every change, that a second, independently written program gets the same answer as the app --- and that both get the right answer on devices whose behaviour is known exactly.
\item[{\hypertarget{term:phantom}{phantom capture}}] A made-up recording whose distortion is known to the digit, used to check that the analysis reads back exactly what was put in.
\item[{\hypertarget{term:phase-removal}{phase removal}}] Using the pedal's own measured harmonic timing to undo what its tone filters did to the picture, so that what is left open is memory and not filtering.
\item[{\hypertarget{term:polarity}{polarity}}] Whether the pedal turns the wave upside down. Stated only when the measurement is sure; otherwise nothing is claimed.
\item[{\hypertarget{term:gate}{presence gate}}] The fixed bar --- four times the noise beside it --- an invented tone must clear to count as measured.
\item[{\hypertarget{term:product-order}{product order} ($|m| + |n|$)}] How many multiples of the two notes go into an invented tone; the app lists every combination up to seven.
\item[{\hypertarget{term:product-reading}{product reading}}] What the app says about each invented tone: measured, nothing found down to a stated level, or could not be read --- and why.
\item[{\hypertarget{term:rate-constant}{rate constant} ($L$)}] How quickly the test tone climbs; a larger value means a slower, longer sweep.
\item[{\hypertarget{term:scatter}{read scatter}}] How far a single measured level might be off simply because of noise, stated alongside the level.
\item[{\hypertarget{term:reduced-ratio}{reduced tone ratio} ($a : b$)}] The two notes' frequencies as a simplest fraction; a fraction with big enough numbers keeps every invented tone on its own slot.
\item[{\hypertarget{term:alignment-error}{residual alignment error}}] A leftover timing error between the test signal and the recording of a fraction of a sample or more; it opens the picture into a loop that is wider for higher notes, and is a property of the rig, not the pedal.
\item[{\hypertarget{term:pair-statistic}{sideband pair statistic}}] The two notes' neighbouring slots compared with each other, a shape that hints at what kind of drift occurred without claiming to know.
\item[{\hypertarget{term:static-curve}{static curve}}] The single line drawn through the middle of a loop, taken as the device's clipping shape: the upward and downward passes averaged together.
\item[{\hypertarget{term:averaging}{sweep averaging}}] Playing the sweep several times in a row and combining the results so that random noise partly cancels while the device's behaviour adds up.
\item[{\hypertarget{term:synchronization}{synchronization condition}}] A small adjustment to the sweep's length so that the arithmetic that separates the harmonics works out exactly.
\item[{\hypertarget{term:sweep}{synchronized swept sine} ($x(t)$)}] A test tone that glides smoothly from low to high pitch, timed so that the distortion each note produces can be picked apart afterwards.
\item[{\hypertarget{term:tanh}{tanh}}] The hyperbolic tangent, the textbook soft clipper: a smooth S-shaped curve that passes small signals through unchanged and bends ever more gently toward a ceiling it never quite reaches. It clips the top and bottom of the wave alike, so it adds only odd harmonics, and they fall away quickly. The worked example in these documents is this curve with the input doubled first.
\item[{\hypertarget{term:thd}{total harmonic distortion} ($\mathrm{THD}$)}] The size of everything the device added at multiples of the note, compared with the note itself, as a single number.
\item[{\hypertarget{term:transfer-curve}{transfer curve}}] A picture of what comes out against what went in, drawn while one steady low note plays: a pedal with no memory draws one line, and a loop means the output depends on where the signal has been as well.
\item[{\hypertarget{term:two-tone}{two-tone measurement}}] The app plays two notes at once --- a power chord's fifth by default --- and measures every extra tone the pedal invents from the pair.
\item[{\hypertarget{term:lobe-area}{unsigned lobe area}}] How much area the loop encloses, counted lobe by lobe regardless of which way each lobe is traced --- a figure-eight counts twice, not zero.
\item[{\hypertarget{term:validity-bound}{validity bound}}] The range of notes over which a given harmonic can be measured honestly; beyond it the recording itself gets in the way.
\end{description}
